\chapter{2013年湖北卷（文）}

\section{单选题}

\begin{problem}
已知全集 \(U = \{1, 2, 3, 4, 5\}\)，集合 \(A = \{1, 2\}\)，\(B = \{2, 3, 4\}\)，则 \(B \cap \complement_U A =\)（\quad）

\choices
    {\( \{2\} \)}
    {\(\{3,4\}\)}
    {\( \{1, 4, 5\} \)}
    {\(\{2, 3, 4, 5\}\)}
\end{problem}

\begin{problem}
已知 \(0 < \theta < \frac{\pi}{4}\)，则双曲线 \(C_{1}: \frac{x^{2}}{\sin^{2}\theta} - \frac{y^{2}}{\cos^{2}\theta} = 1\) 与 \(C_{2}: \frac{y^{2}}{\cos^{2}\theta} - \frac{x^{2}}{\sin^{2}\theta} = 1\) 的（\quad）

\choices
    {实轴长相等}
    {虚轴长相等}
    {离心率相等}
    {焦距相等}
\end{problem}

\begin{problem}
在一次跳伞训练中，甲、乙两位学员各跳一次. 设命题 \(p\) 是“甲降落点在指定范围”，\(q\) 是“乙降落在指定范围”，则命题“至少有一位学员没有降落在指定范围”可表示为（\quad）

\choices
    {\((-p)\vee (-q)\)}
    {\(p\vee (-q)\)}
    {\((-p)\wedge (-q)\)}
    {\(p\vee q\)}
\end{problem}

\begin{problem}
四名同学根据各自的样本数据研究变量 \(x, y\) 之间的相关关系，并求得回归直线方程，分别得到以下四个结论： \circled{1} y 与 x 负相关且 \( \beta = 2.347x - 6.423 \) ; \circled{2} y 与 x 负相关且 \( \beta = -3.470x + 5.648 \) ; \circled{3} y 与 x 正相关且 \( \beta = 5.437x + 8.493 \) ; \circled{4} y 与 x 正相关且 \( \beta = -4.320x - 4.578 \) . 其中一定不正确的结论的序号是（\quad）

\choices
    {\circled{1}\circled{2}}
    {\circled{2}\circled{3}}
    {\circled{3}\circled{4}}
    {\circled{1}\circled{4}}
\end{problem}

\begin{problem}
小明骑车上学，开始均匀速行驶，途中因交通堵塞停留了一段时间后，为了赶时间加快速度行驶，与以上事件吻合得最好的图象是（\quad）

\choices
    {\bitmapinclude[max width=3.3cm,max totalheight=2.8cm]{img/2013/hubei_liberal/q5_choice_a.png}}
    {\bitmapinclude[max width=3.3cm,max totalheight=2.8cm]{img/2013/hubei_liberal/q5_choice_b.png}}
    {\bitmapinclude[max width=3.3cm,max totalheight=2.8cm]{img/2013/hubei_liberal/q5_choice_c.png}}
    {\bitmapinclude[max width=3.3cm,max totalheight=2.8cm]{img/2013/hubei_liberal/q5_choice_d.png}}
\end{problem}

\begin{problem}
将函数 \( y = \sqrt{3}\cos x + \sin x (x \in \R) \) 的图象向左平移 \( m (m > 0) \) 个单位长度后，所得到的图象关于 \( y \) 轴对称，则 \( m \) 的最小值是（\quad）

\choices
    {\(\frac{3}{12}\)}
    {\(\frac{\pi}{6}\)}
    {\(\frac{\pi}{3}\)}
    {\(\frac{5\pi}{6}\)}
\end{problem}

\begin{problem}
已知点 \(A(-1,1), B(1,2), C(-2,-1), D(3,4)\)，则向量 \(\overrightarrow{AB}\) 在 \(\overrightarrow{CD}\) 方向上的投影为（\quad）

\choices
    {\(\frac{3\sqrt{2}}{2}\)}
    {\(\frac{3\sqrt{15}}{2}\)}
    {\(-\frac{3\sqrt{2}}{2}\)}
    {\(-\frac{3\sqrt{15}}{2}\)}
\end{problem}

\begin{problem}
\(x\) 为实数， \([x]\) 表示不超过 \(x\) 的最大整数，则函数 \(f(x) = x - [x]\) 在 \(\R\) 上为（\quad）

\choices
    {奇函数}
    {偶函数}
    {增函数}
    {周期函数}
\end{problem}

\begin{problem}
某旅行社租用 A, B 两种型号的客车安排 900 名客人旅行，A, B 两种车辆的载客量分别为 36 h 和 60 h, A 租金分别为 1600 元/辆和 2400 元/辆. 旅行社要求租车总数不超过 21 辆. 且 B 型车不多于 A 型车 7 辆，则租金最少为（\quad）

\choices
    {31200 元}
    {36000 元}
    {36800 元}
    {38400 元}
\end{problem}

\begin{problem}
已知函数 \( f(x) = x (\ln x - ax) \) 有两个极值点，则实数 \( a \) 的取值范围是（\quad）

\choices
    {\( (-\infty,0) \)}
    {\( \left(0,\frac{1}{2}\right) \)}
    {\( (0,1) \)}
    {\( (0,+\infty) \)}
\end{problem}

\section{填空题}

\begin{problem}
i 为虚数单位，设复数 \( z_{1} \) ， \( z_{2} \) 在复平面内对应的点关于原点对称，若 \( z_{1}=2-3\i \) ，则 \( z_{2}= \)  \fillinblank{}.
\end{problem}

\begin{problem}
某学员在一次射击测试中射靶 10 次，命中环数如下: 7, 8, 7, 9, 5, 4, 9, 10, 7, 4 则：(1) 平均命中环数为 \fillinblank{}；(2) 命中环数的标准差为 \fillinblank{} \fillinblank{}.
\end{problem}

\begin{problem}
阅读如图所示的程序框图，运行相应的程序，若输入 m 的值为 2，则输出的结果是 \fillinblank{} \fillinblank{}.
\end{problem}

\begin{problem}
已知圆 \(O: x^{2} + y^{2} = 5\)，直线 \(l: x \cos \theta + y \sin \theta = 1 (0 < \theta < \frac{\pi}{2})\). 设圆 \(O\) 上到直线 \(l\) 的距离等于 1 的点的个数为 \(k\)，则 \(k = \)  \fillinblank{}.
\end{problem}

\begin{problem}
在区间 \([-2, 4]\) 上随机地取一个数 \(x\)，若 \(x\) 满足 \(|x| \leqslant m\) 的概率为 \(\frac{5}{6}\)，则 \(m = \)  \fillinblank{}.
\end{problem}

\begin{problem}
我国古代数学名著《数学九章》中有“天池盈渊则”题. 在下河时，用一个圆台形的大池盆接雨水，天池盆盆口直径为二尺八寸，盆底直径为一尺二寸，盆深一尺八寸，若盆中积水深九寸，则平地降雨量是寸. (注：\circled{1} 平地降雨量等于盆中积水体积除以盆口面积；\circled{2} 一尺等于十寸)
\end{problem}

\begin{problem}
在平面直角坐标系中，若点 \(P(x, y)\) 的坐标为 \(x\)，\(y\) 均为整数，则称点 \(P\) 为格点，若一个多边形的顶点全是格点，则称该多边形为格点多边形，格点多边形的面积为 \(S\)，其内部的格点数记为 \(N\)，边界上的格点数记为 \(L\). 例如图中 \(\triangle ABC\) 是格点三角形，对应的 \(S = 1, N = 0, L = 4\).

(1) 图中格点四边形 DEFG 对应的 S, N, L 分别是 ; (2) 已知格点多边形的面积可表示为 \( S = aN + bL + c \) ，其中 a, b, c 为常数. 若某格点多边形对应的 N = 71, L = 18，则 S =\fillinblank{}. （用数值作答）
\end{problem}

\section{解答题}

\begin{problem}
在 \(\triangle ABC\) 中，角 \(A, B, C\) 对应的边分别是 \(a, b, c\)，已知 \(\cos 2A - 3\cos (B + C) = 1\). (1) 求角 A 的大小; (2) 若 \( \triangle ABC \) 的面积 \( S = 5\sqrt{3} \) ，b = 5，求 \( \sin B \sin C \) 的值.
\end{problem}

\begin{problem}
已知 \(S_{n}\) 是等比数列 \(\{a_{n}\}\) 的前 \(n\) 项和，\(S_{4}, S_{2}, S_{3}\) 成等差数列. 且 \(a_{2} + a_{3} + a_{4} = -18\). (1) 求数列 \( \{a_{n}\} \) 的通项公式; (2) 是否存在正整数 \(n\)，使得 \(S_{n} \geqslant 2013\)\(\perp\) 若存在，求出符合条件的所有 \(n\) 的集合; 若不存在，说明理由.
\end{problem}

\begin{problem}
设 \(a > 0, b > 0\)，已知函数 \(f(x) = \frac{ax + b}{a + b}\). (1) 当 \(a \neq b\) 时，讨论函数 \(f(x)\) 的单调性. (2) 当 \(x > 0\) 时，称 \(f(x)\) 为 \(a, b\) 关于 \(x\) 的加权平均数. \circled{1} 判断 \( f(1), f\left(\sqrt{\frac{b}{a}}\right), f\left(\frac{b}{a}\right) \) 是否成等比数列，并证明 \( f\left(\frac{b}{a}\right) \leqslant f\left(\sqrt{\frac{b}{a}}\right) \); \circled{2} a、b 的几何平均数记为 G，称 \( \frac{2ab}{a+b} \) 为 a, b 的调和平均数，记为 H，若 \( H \leqslant f(x) \leqslant G \) ，求 x 的取值范围.
\end{problem}

\begin{problem}
如图，已知椭圆 \(C_1\) 与 \(C_2\) 的中心在坐标原点 \(O\)，长轴均为 \(MN\) 且在 \(x\) 轴上，短轴长分别为 \(2m, 2n (m > n)\)，过原点且不与 \(x\) 轴重合的直线 \(l\) 与 \(C_1, C_2\) 的四个交点按纵坐标从大到小依次为 \(A, B, C, D\)，记 \(\lambda = \frac{1}{m}, \triangle BDM\) 和 \(\triangle ABN\) 的面积分别为 \(S_1\) 和 \(S_2\). (1) 当直线 l 与 y 轴重合时，若 \( S_{1}=AS_{2} \) ，求 \( \lambda \) 的值; (2) 当 \(\lambda\) 变化时，是否存在与坐标轴不重合的直线 \(l\)，使得 \(S_{1} = \lambda S_{2}\)\(\perp\) 并说明理由.
\end{problem}

\begin{problem}
如图，某地旅队自水平地面 \(A, B, C\) 三处垂直向地下钻探，自 \(A\) 点向下钻到 \(A_{1}\) 处发现矿藏，再继续下钻到 \(A_{2}\) 处后下面已无矿，从而得到在 \(A\) 处正下方的矿层厚度为 \(A_{1}A_{2} = d_{1}\)，同样可得在 \(B, C\) 处正下方的矿层厚度分别为 \(B_{1}B_{2} = d_{2}, C_{1}C_{2} = d_{3}\)，且 \(d_{1} < d_{2} < d_{3}\)，过 \(AB, AC\) 的中点 \(M, N\) 且与直线 \(AA_{2}\) 平行的平面数多面体 \(A_{1}B_{1}C_{1} - A_{2}B_{2}C_{2}\) 所得的截面 \(DEFG\) 为该多面体的一个截面，其面积记为 \(S_{\mathrm{D}}\). (1) 证明: 中截面 \(DEFG\) 是梯形. (2) 在 \(\triangle ABC\) 中，已知 \(BC = a\)，\(BC\) 边上的高为 \(h\)，面积为 \(S\). 在估测三角形 \(ABC\) 区域内正下方的矿囊储量 (即多面体 \(A_{1}B_{1}C_{1} - A_{2}B_{2}C_{2}\) 的体积 \(V\)) 时，可用近似公式 \(V_{\text{内}} = S_{\text{中}} \cdot h\) 来估算. 已知 \(V = \frac{1}{3}(d_{1} + d_{2} + d_{3})S\)，试判断 \(V_{\text{内}}\) 与 \(V\) 的大小关系，并加以证明.
\end{problem}

\begin{problem}
如图，已知椭圆 \(C_1\) 与 \(C_2\) 的中心在坐标原点 \(O\)，长轴均为 \(MN\) 且在 \(x\) 轴上，短轴长分别为 \(2m, 2n (m > n)\)，过原点且不与 \(x\) 轴重合的直线 \(l\) 与 \(C_1, C_2\) 的四个交点按纵坐标从大到小依次为 \(A, B, C, D\)，记 \(\lambda = \frac{m}{n}, \triangle BDM\) 和 \(\triangle ABN\) 的面积分别为 \(S_1\) 和 \(S_2\). (1) 当直线 l 与 y 轴重合时，若 \( S_{1}=AS_{2} \) ，求 \( \lambda \) 的值; (2) 当 \(\lambda\) 变化时，是否存在与坐标轴不重合的直线 \(l\)，使得 \(S_{1} = \lambda S_{2}\)\(\perp\) 并说明理由.
\end{problem}

